% Fault101 Fault 2_3 Explorer - Step-by-Step Student Guide
% Generated by guide_tools/latex/latexkit.py from the story script; this file is plain LaTeX - edit it directly.
% Build: pdflatex <this file>   (run it twice so the contents page fills in).  Needs scaguide.sty and figures/ next to it.
\documentclass[11pt,letterpaper]{article}
\usepackage{scaguide}
\guidesetup{
  headerleft={ChipWhisperer Fault101  -  Fault 2\_3 Explorer Guide},
  headerright={Memory Dump Explorer},
  footerleft={Companion to: Fault 2\_3 - Voltage Glitching to Memory Dump Explorer.ipynb},
  title={Fault101 Fault 2\_3 Explorer - Step-by-Step Student Guide},
  subject={Step-by-step guide to the Fault 2\_3 Voltage Glitching Explorer: a response-length bug that is easy to glitch, in a direction the lab's name does not predict}}

\begin{document}

\guidecover{FAULT101  -  FAULT 2\_3  (EXPLORER)}{Memory Dump\\ Explorer}{Step-by-step student guide: a bug that is easy to glitch - but maybe not the way you expect}

Fault 2\_3 targets a bootloader's reply routine, whose length is controlled by a variable a fault can corrupt. The
lab's name promises a \textbf{memory dump} - a fault that makes the device echo back more than it should. This guide's
central finding is that on our kit, \textbf{the fault goes the other way}: it reliably shortens the reply instead of
lengthening it. Both directions come from the exact same root cause, this one is just much easier to find than Fault
2\_2's password bypass, and - like Fault 2\_2 - the exact settings that work \textbf{drift} over the course of a
session. \textbf{This lab needs real hardware; there is no simulated mode.}

\vspace{10pt}

\begin{guidecallout}{tip}{HOW TO USE THIS GUIDE}
This guide goes \textbf{with} the notebook - it does not replace it. Keep the notebook open next to this PDF.
\end{guidecallout}

\guidehthree{\textbf{Which files are which}}

{\small\arrayrulecolor{guidegrid}
\rowcolors{2}{white}{guidezebra}
\begin{longtable}{|R{0.4348}|R{0.4493}|R{0.1159}|}
\hline
\rowcolor{guidenavy}\hdrcell{File (folder: jupyter/courses/fault101/)} & \hdrcell{What it is} & \hdrcell{Open it?} \\ \hline
\endhead
\codesmall{Fault 2\_\bk{}3 - Voltage Glitching to Memory Dump Explorer.\bk{}ipynb} & The Explorer. Everything in this guide happens here. & \textbf{Yes} \\ \hline
\codesmall{Fault 2\_\bk{}3 - Voltage Glitching to Memory Dump.\bk{}ipynb} & The original (non-Explorer) lab. & Optional \\ \hline
\codesmall{SOLN\_\bk{}Fault 2\_\bk{}3 - .\bk{}.\bk{}.\bk{}ipynb} & Answer key for the original lab. & If stuck \\ \hline
\codesmall{firmware/\bk{}mcu/\bk{}bootloader-glitch/\bk{}bootloader.\bk{}c} & (a different folder: \codesmall{chipwhisperer/\bk{}firmware/\bk{}.\bk{}.\bk{}.\bk{}}) The reply routine this lab targets. & Background \\ \hline
\end{longtable}}

\vspace{4pt}

\guidehthree{\textbf{The four colored boxes you will see everywhere}}

\guidelegend

\vspace{4pt}

\guidesmall{Guide version 1.0 - November 2026. Every number in this guide was measured on one real kit:
ChipWhisperer-Husky, a CW313 board and a CW312T-SAM4S target (Windows 11, Python 3.12, ChipWhisperer 6.0.0, firmware
1.5.0) - the same kit as the Fault 2\_1 and 2\_2 guides. As with Fault 2\_2, \textbf{the exact settings below drift
within a single session} - treat specific numbers as `this was possible', not as guarantees.}

\clearpage

\tableofcontents

\clearpage

\guidehone{The big picture}

\textbf{What is the bug?} The test bootloader receives a hex-encoded upload, `decrypts' it (a ROT-13 substitution -
nothing secret), and replies with a fixed acknowledgement. A loop sends back \code{ascii\_idx} bytes from a buffer,
and nothing re-checks \code{ascii\_idx} between being set and being used in that loop:

\begin{guidecode}
ascii_buffer[ascii_idx++] = 'r'; ascii_buffer[ascii_idx++] = '0';   /* + 6x '\n' */
...
for(i = 0; i < ascii_idx; i++)
    putch(ascii_buffer[i]);
\end{guidecode}

If a fault corrupts \code{ascii\_idx} (or the loop's comparison against it) in between, the reply is the wrong length.
Corrupt it \textbf{upward} past what was written and you leak extra bytes sitting after it in memory - the over-read
this lab is named for. Corrupt it \textbf{downward} and the reply is silently \textbf{truncated} instead - still a
real bug (a caller expecting a fixed-length reply gets less and may misread what is missing), just a quieter one.
Which direction you get depends on exactly how the compiler laid things out - which is exactly the point of this
Explorer: it has you find out which direction happens on your own build, rather than assuming the dramatic one.

\guidehthree{Your path through the notebook:}

\begin{center}
\begin{tikzpicture}
\node[rounded corners=4pt,fill=guidenavy,text=white,align=center,minimum height=1.15cm,minimum width=3.96cm,inner sep=3pt] (n0) at (1.98,0) {\textbf{1  Baseline}\\{\footnotesize and reuse Fault 2\_2}};
\node[rounded corners=4pt,fill=guidenavy,text=white,align=center,minimum height=1.15cm,minimum width=3.96cm,inner sep=3pt] (n1) at (6.49,0) {\textbf{2 - 3  Map}\\{\footnotesize the crash edge}};
\node[rounded corners=4pt,fill=guidenavy,text=white,align=center,minimum height=1.15cm,minimum width=3.96cm,inner sep=3pt] (n2) at (11.01,0) {\textbf{4  Check}\\{\footnotesize reliability + payload}};
\node[rounded corners=4pt,fill=guideteal,text=white,align=center,minimum height=1.15cm,minimum width=3.96cm,inner sep=3pt] (n3) at (15.52,0) {\textbf{5  Which way}\\{\footnotesize does it fail?}};
\draw[-{Stealth},guidegrey,thick] (n0.east) -- (n1.west);
\draw[-{Stealth},guidegrey,thick] (n1.east) -- (n2.west);
\draw[-{Stealth},guidegrey,thick] (n2.east) -- (n3.west);
\end{tikzpicture}
\end{center}

\vspace{6pt}

{\small\arrayrulecolor{guidegrid}
\rowcolors{2}{white}{guidezebra}
\begin{longtable}{|R{0.1159}|R{0.4348}|R{0.4493}|}
\hline
\rowcolor{guidenavy}\hdrcell{Experiment} & \hdrcell{What you do} & \hdrcell{What you learn} \\ \hline
\endhead
\textbf{1} & Try Fault 2\_1/2\_2's settings here too. & Still nothing - the interesting part of this routine's window is elsewhere. \\ \hline
\textbf{2} & Scan width at a later moment to find the crash/success edge. & A clean, easy-to-find edge: past it, almost every shot truncates the reply - at that instant. \\ \hline
\textbf{3} & Scan a wide range of moments at a reliable width. & Several good moments scattered through the window - far denser than Fault 2\_2's. \\ \hline
\textbf{4} & Check reliability and whether the payload matters. & The same length truncates every time at a given setting, regardless of what you uploaded. \\ \hline
\textbf{5} & Search specifically for a longer-than-baseline reply. & On our kit, never found - only ever truncation, across thousands of attempts. \\ \hline
\end{longtable}}

\clearpage

\guidehone{Before you start}

Same kit and checklist as Fault 2\_1 and 2\_2.

{\small\arrayrulecolor{guidegrid}
\rowcolors{2}{white}{guidezebra}
\begin{longtable}{|R{0.0652}|R{0.4783}|R{0.4565}|}
\hline
\rowcolor{guidenavy}\hdrcell{} & \hdrcell{You need} & \hdrcell{How to check} \\ \hline
\endhead
[  ] & Fault 2\_1 and 2\_2 completed (same setup and hardware). & You can fire a glitch and read the outcomes. \\ \hline
[  ] & \textbf{ChipWhisperer-Husky} plugged straight into a USB port. & Lab 0's hardware test ends with \codesmall{hardware OK (disconnected again)}. \\ \hline
[  ] & \textbf{CW312T-SAM4S} on \textbf{CW313}, ribbon connected. & See the SCA101 Linux guide, Part 4.2. \\ \hline
[  ] & The Husky \textbf{\codesmall{Crowbar}} SMA cable connected. & Without it, every glitch here shows as \codesmall{normal} too. \\ \hline
[  ] & No other notebook holding the Husky open. & \textbf{File \textgreater{} Shut Down Notebook} elsewhere first. \\ \hline
\end{longtable}}

\vspace{4pt}

\begin{guidecallout}{tip}{GOOD TO KNOW}
\textbf{This one is easier than Fault 2\_2.} A plain width-then-moment scan found working settings in well under a
minute on our kit - no blind random search over a huge box was needed. Do not let Fault 2\_2's difficulty set your
expectations for every lab.
\end{guidecallout}

\guidehone{Experiment 1 - try Fault 2\_1/2\_2's settings here too}

\textbf{Goal:} check, rather than assume, whether a setting that worked for the counting loop or the password check
does anything to this response routine.

\begin{guidecallout}{see}{YOU SHOULD SEE}
All \code{normal}. \code{ext\_offset=1} is early in this firmware's own (much longer) trigger window - which starts
when the upload begins, not when the reply is sent - and this width is too gentle for this firmware's own crash
threshold anyway. The interesting part of the window is much later.
\end{guidecallout}

\guidehone{Experiment 2 - find this firmware's own crash edge}

\textbf{Goal:} scan width at a fixed offset, somewhere later in the trigger window than Experiment 1
(\code{ext\_offset} around two-thirds of the way through \code{trig\_count}), the way Fault 2\_1 and 2\_2 did for
their own crash edges.

\begin{guidecallout}{see}{YOU SHOULD SEE}
At one particular \code{(offset, ext\_offset)}: width 1775 and below was \code{normal}; 1800-1825 a transition; and
\textbf{1850 and above gave \code{success} every single time}, up to 2125, every reply shortened to the same 6 bytes.
That is how clean this edge can look at one instant.
\end{guidecallout}

\begin{guidecallout}{warn}{GOOD TO KNOW}
\textbf{This moves too, like Fault 2\_2's settings.} Re-running the exact same cell roughly twenty minutes later in
the same session, that same \code{ext\_offset} had gone completely quiet - \code{normal} at every width from 1600 to
2125. Experiment 3 found a different \code{ext\_offset} nearby that had picked up the effect instead, and the same
kind of edge reappeared there, just less sharp (a mix of outcomes above about 1825, not a clean 100\%). \textbf{Don't
treat any single \code{(width, offset, ext\_\bk{}offset)} as permanent} - if a setting that worked goes quiet, check
its neighbours (especially in \code{ext\_offset}) before assuming the effect disappeared.
\end{guidecallout}

\guidehone{Experiment 3 - how much of the window is like this?}

\textbf{Goal:} Experiment 2 picked one moment almost at random. Scan a reliably-strong width across the *whole*
trigger window to see how much of it behaves the same way.

\begin{guidecallout}{see}{YOU SHOULD SEE}
Mostly \code{normal}, with \code{success} at a handful of specific moments scattered through the window - in one run,
around 750, 1350, 1800 and 2100 out of roughly 2150 total; in another, 1500, 1800 and 2100. Roughly one good moment
every few hundred cycles - far denser than Fault 2\_2 ever found for its password bypass. Here, a plain scan over
width and a handful of moments, with no blind random search, turns up several working settings in well under a minute.
\end{guidecallout}

\guidehone{Experiment 4 - reliability and payload-independence}

\textbf{Goal:} confirm a currently-working setting is reproducible, and check whether the uploaded message itself
makes any difference - this fault corrupts a loop bound, not the message's content, so it should not.

\begin{guidecallout}{see}{YOU SHOULD SEE}
Two very different payloads (the original hex sentence, and a message of all \code{0x41} bytes) behaved almost
identically at a setting that was working at the time: 14 of 20 and 12 of 20 shots succeeded - close enough to call
`the same rate', given how much these settings drift on their own (Experiment 2). \textbf{Every success truncated to
the same length} (6 bytes) regardless of which payload was sent. That confirms the fault corrupts a loop bound or
counter, not the content you uploaded. Different settings truncated to different lengths across our testing (we also
saw replies of 1, 2, 3, 4, 5 and 7 bytes), but a *given* setting kept giving the *same* length for as long as it kept
working at all.
\end{guidecallout}

\guidehone{Experiment 5 - which direction does it fail in, and why does that matter?}

\textbf{Goal:} this lab is named `memory dump' because the intended bug is an \textbf{over-read} - a fault making
\code{ascii\_idx} too big, so the reply leaks bytes that were never meant to go out. Search specifically for a reply
longer than the 8-byte baseline, across a wide box of settings.

\begin{guidecallout}{see}{YOU SHOULD SEE}
Across every search we ran on this kit - several thousand attempts in total, spread over a wide box of widths, offsets
and moments - \textbf{the reply was never longer than the 8-byte baseline.} The fault only ever truncated it. That is
the real finding here: the same root cause produced an under-read, not the over-read the lab's name promises.
\end{guidecallout}

\begin{guidecallout}{tip}{GOOD TO KNOW}
\textbf{Neither direction is `better' or `worse' to report.} A silent truncation can be just as damaging in a protocol
where a caller trusts a length field and a short reply means something was quietly dropped. Recording which direction
actually happened, honestly, is more useful than assuming the dramatic one did. If your own kit or a different
compiler build ever produces a longer-than-baseline reply, that is the over-read this lab was named for - write down
the exact setting and what extra bytes came back.
\end{guidecallout}

\guidehone{Challenges (optional)}

\begin{enumerate}[label=\arabic*.,leftmargin=1.8em,itemsep=2pt,topsep=2pt,parsep=0pt,start=1]
\item \textbf{Map the full truncation spectrum.} Collect every distinct reply length you can find across several settings. Is there a pattern to which settings give which length?
\item \textbf{Compare search difficulty properly.} Time a plain width-then-moment scan here against Fault 2\_2's blind random search. Which bug was easier to glitch, and by roughly how much?
\item \textbf{Does \code{repeat} change the truncation length?} Try \code{REPEAT = 2} or \code{3} at a working setting.
\item \textbf{A longer payload.} Upload a message close to the buffer's limit (80 hex characters) and repeat Experiment 2. Does the crash edge or the truncated length change?
\end{enumerate}

\guidehone{If something goes wrong}

{\small\arrayrulecolor{guidegrid}
\rowcolors{2}{white}{guidezebra}
\begin{longtable}{|R{0.2899}|R{0.3043}|R{0.4058}|}
\hline
\rowcolor{guidenavy}\hdrcell{What you see} & \hdrcell{Most likely cause} & \hdrcell{What to do} \\ \hline
\endhead
The baseline assert fails. & Firmware did not flash, or a stale serial buffer. & Re-run the build/flash cells; re-run \codesmall{reboot\_flush()}. \\ \hline
Every shot is \codesmall{normal}, even at a huge width. & The Crowbar SMA cable is not connected. & Fit the cable, run the setup cells again. \\ \hline
Every shot is \codesmall{reset}. & You are above the crash threshold for this width. & Lower \codesmall{WIDTH} a little. \\ \hline
\codesmall{target.read()} returns nothing at all. & The target may be stuck mid-reply. & Run \codesmall{recover()}, then \codesmall{reboot\_flush()}. \\ \hline
A setting that worked earlier stops working. & Drift (see the note after Experiment 2). & Re-measure with fresh shots; try nearby \codesmall{ext\_offset} values. \\ \hline
\codesmall{OSError: Could not find ChipWhisperer}. & Husky unplugged, or another notebook still holds it. & Plug it in; shut down every other notebook; restart this kernel. \\ \hline
\codesmall{USBErrorPipe} / \codesmall{Unknown hwInfoVer}. & The Husky's USB link or FPGA stopped answering. & Unplug and replug the Husky directly; restart the kernel. \\ \hline
A cell was interrupted (Stop button). & Target or glitch hardware may be left in an odd state. & Run \codesmall{recover()} in a new cell. \\ \hline
\end{longtable}}

\vspace{4pt}

\guidehone{Appendix - every variable you can edit}

{\small\arrayrulecolor{guidegrid}
\rowcolors{2}{white}{guidezebra}
\begin{longtable}{|R{0.2174}|R{0.1449}|R{0.6377}|}
\hline
\rowcolor{guidenavy}\hdrcell{Variable} & \hdrcell{Default} & \hdrcell{What it does} \\ \hline
\endhead
\codesmall{WIDTH}, \codesmall{OFFSET}, \codesmall{EXT\_OFFSET} & varies & The glitch settings fired, in Husky phase-shift steps / clock cycles. \\ \hline
\codesmall{REPEAT} & 1 & Number of glitch pulses fired per shot. \\ \hline
\codesmall{SHOTS} & varies & How many glitches to fire for one measurement. \\ \hline
\codesmall{WIDTH\_\bk{}START/\bk{}STOP/\bk{}STEP} & varies & The width range scanned in Experiment 2. \\ \hline
\codesmall{EXT\_STEP} & 150 & The step used when scanning the whole trigger window in Experiment 3. \\ \hline
\codesmall{PAYLOAD\_A}, \codesmall{PAYLOAD\_B} & two messages & The two uploads compared in Experiment 4. \\ \hline
\codesmall{WIDTH\_RANGE}, \codesmall{OFFSET\_RANGE}, \codesmall{TIME\_BUDGET\_S} & wide / 120 & The box and time budget searched in Experiment 5. \\ \hline
\end{longtable}}

\vspace{4pt}

\guidehone{Appendix - Glossary}

{\small\arrayrulecolor{guidegrid}
\rowcolors{2}{white}{guidezebra}
\begin{longtable}{|R{0.2464}|R{0.7536}|}
\hline
\rowcolor{guidenavy}\hdrcell{Term} & \hdrcell{Meaning} \\ \hline
\endhead
Voltage glitch / fault injection / drift & See the Fault 2\_1 and 2\_2 guides' glossaries - the same concepts apply here. \\ \hline
Over-read & A fault (or bug) that makes code send back more data than intended, potentially leaking memory it should not. \\ \hline
Truncation / under-read & A fault (or bug) that makes code send back less data than intended - the direction this lab's firmware actually took on our kit. \\ \hline
\codesmall{ascii\_idx} & The variable controlling how many bytes the bootloader's reply loop sends back. \\ \hline
\codesmall{trig\_count} & How many ADC samples the trigger stayed high for during one capture - used to size the \codesmall{ext\_offset} search range. \\ \hline
\end{longtable}}

\vspace{4pt}

\guidehthree{Notice}

\guidesmall{This guide accompanies ChipWhisperer course material that references (commercial) training material on
ChipWhisperer.io. You can freely execute and use the lab per the open-source license (including in your own courses if
you distribute similarly), but you must maintain notice about the source: https://www.ChipWhisperer.io.}

\guidesmall{Copyright (C) NewAE Technology Inc., 2015-2020. ChipWhisperer is a trademark of NewAE Technology Inc.,
claimed in all jurisdictions, and registered in at least the United States of America, European Union, and Peoples
Republic of China. Tutorials derived from our open-source work must be released under the associated open-source
license, and notice of the source must be clearly displayed.}

\guidesmall{THE SOFTWARE IS PROVIDED ``AS IS'', WITHOUT WARRANTY OF ANY KIND, EXPRESS OR IMPLIED, INCLUDING BUT NOT
LIMITED TO THE WARRANTIES OF MERCHANTABILITY, FITNESS FOR A PARTICULAR PURPOSE AND NONINFRINGEMENT. IN NO EVENT SHALL
THE AUTHORS OR COPYRIGHT HOLDERS BE LIABLE FOR ANY CLAIM, DAMAGES OR OTHER LIABILITY, WHETHER IN AN ACTION OF
CONTRACT, TORT OR OTHERWISE, ARISING FROM, OUT OF OR IN CONNECTION WITH THE SOFTWARE OR THE USE OR OTHER DEALINGS IN
THE SOFTWARE.}

\end{document}
